(a) The germ $m_{\mathfrak p}$ vanishes exactly when some $a\notin\mathfrak p$ annihilates $m$. Thus $m_{\mathfrak p}\ne0$ exactly when $\operatorname{Ann}(m)\subseteq\mathfrak p$, giving $\operatorname{Supp}(m)=V(\operatorname{Ann}(m))$.

(b) Choose generators $m_1,\ldots,m_r$ of $M$. If $M_{\mathfrak p}=0$, choose $a_i\notin\mathfrak p$ with $a_im_i=0$. Then $\prod_i a_i\in\operatorname{Ann}(M)\setminus\mathfrak p$. The converse follows by localizing an annihilator outside $\mathfrak p$. Hence $\operatorname{Supp}(\widetilde M)=V(\operatorname{Ann}(M))$.

(c) Apply (b) on an affine open covering.

(d) For $j:X\setminus Z\hookrightarrow X$, the sheaf $\mathcal H_Z^0(\widetilde M)=\ker(\widetilde M\to j_*j^*\widetilde M)$ is quasi-coherent, since $X$ is noetherian. By (a), its global sections are the elements $m$ satisfying $\mathfrak a\subseteq\sqrt{\operatorname{Ann}(m)}$. Since $\mathfrak a$ is finitely generated, this is equivalent to $\mathfrak a^nm=0$ for some $n$. Its module of global sections is therefore $\Gamma_{\mathfrak a}(M)$, which gives the required isomorphism.

(e) Quasi-coherence follows locally from (d). If $\mathcal F$ is coherent, $\mathcal H_Z^0(\mathcal F)$ is locally associated to a submodule of a finite module over a noetherian ring, hence is coherent.
